\addsectionframe{7. Ausblick und Fazit}

\begin{itemframe}{Erweiterbarkeit konkret}
  \item Die Anleitung \texttt{docs/processes/architecture.md} beschreibt das Hinzufügen von Modulen und Themen.
  \item Dazu gehören Ordnerstruktur, gemeinsame Seitengerüste, Registrierung im \texttt{moduleCatalog} und Übersetzungen.
  \item Offene Lücke: Diese Anleitung sollte vor Abgabe mit einem neuen Beispielthema gegengeprüft werden.
\end{itemframe}

\begin{itemframe}{Nächste sinnvolle Schritte}
  \item Echtes Feedback aus Lehrveranstaltungen dokumentieren und in Anforderungen zurückführen.
  \item Produktiv-Deployment mit Hosting, Freigabe und Rückfallplan festlegen.
  \item Gamification, LMS-Anbindung und weitere mobile Optimierungen erst auf dieser Basis priorisieren.
\end{itemframe}

\begin{itemframe}{Fazit}
  \item Die Plattform verbindet fachliche Modelle mit bearbeitbaren, schrittweisen Simulationen.
  \item Architektur und gemeinsame UI-Bausteine halten die Fachmodule konsistent und erweiterbar.
  \item Die Live-Demo prüft den Kernanspruch unmittelbar: Die Plattform erklärt sich durch Benutzung.
\end{itemframe}
